%========================
% Prologue — The Speculative Vessel (Proposal)
%========================

% \chapter*{Prologue — The Speculative Vessel} % or \chapter if numbered in your class

This prologue frames the dissertation as a passage by vessel. The “vessel” names both instrument and ship, both an interface and a transportational force. Here, sound is treated as a navigational medium: it steers perception, organizes bodies, and encodes memory. The prologue sets the tone and the stance: tools are never neutral; they are built things with routes and rules. The work that follows proposes \textit{Speculocultural Technopoiesis} (ST) as a way to redesign those routes from the perspectives of Black sonic epistemologies.

\subsubsection{Objectives}
\begin{itemize}[nosep]
	\item Establish the vessel metaphor as a concise frame for mediation and method.
	\item Declare positionality and ethical commitments (specificity, collaboration, transmission).
	\item Provide a brief reading guide that orients the reader to the structure and stakes.
\end{itemize}

\subsubsection{Intended Sections/Subsections}
\begin{itemize}[nosep]
	\item Purpose and Tone
	\item Researcher Positionality
	\item Orienting Metaphor
\end{itemize}
%-----------------------------------------
\begin{comment}
\subsubsection{Purpose and Tone}
The purpose is to welcome the reader into a method that couples speculation with making. The tone is invitational and exacting: poetic enough to mark what is at stake, precise enough to guide rigorous practice.

\subsubsection{Positionality and Commitments}
The researcher is positioned as both maker and witness. Commitments include cultural specificity over “universal” abstraction; collaborative authorship; documentation that couples code with rationale; and practices of citation that honor lineages.

\subsubsection{Orienting Metaphor and Reading Guide}
\textbf{Vessel as interface.} Interfaces are hulls: they keep some waters out, carry some cargo in, and determine what voyages are possible.  
\textbf{How to read.} Chapter~1 unsettles neutrality; Chapter~2 assembles theory; Chapter~3 specifies method; Chapter~4 demonstrates via case studies; Chapter~5 interprets; Chapter~6 consolidates and looks forward.


% --- Prologue Key References (keyword-based filter) ---
\begin{refsection}[references.bib] % ensure items you want here have keyword=prologue
	\nocite{*}
	\printbibliography[keyword=prologue,heading=subbibliography,title={Key References}]
\end{refsection}
\end{comment}

